\section{Conclusion}
\label{section:conclusion}

We built an end-to-end system that transcribes printed monophonic staff notation
directly into Jianpu, trained entirely on a controlled synthetic corpus whose
generator emits four decoupled label streams as exact ground truth. The
decoupled-stream design turned an architecture comparison into a diagnosis: a
pretrained ViT encoder with an autoregressive decoder solves the task
($\mathrm{SER}=0.0029$), a from-scratch ResNet encoder appears to fail at pitch
under the same decoder, yet a CTC decoder over the \emph{identical} ResNet
features attains a perfect $\mathrm{SER}=0.0000$---locating the bottleneck in the
decoder's monotonic alignment rather than the encoder. The same CTC model is also
far more sample-efficient, reaching usable quality from $5{,}000$ samples where
ViT-AR needs $50{,}000$, and a targeted photo-augmentation fine-tune more than
halves error on real photographs, with geometric (camera-angle) augmentation the
dominant lever. Together these results argue that, for music notation, the
decoder's alignment inductive bias---not the choice of pretrained visual
encoder---is what governs whether and how cheaply pitch is learned.
